% ======================================================================
% SECTION 6: CONCLUSIONS
% Summary of accomplishments and persistent themes. 4 to 6 paragraphs.
% ======================================================================

[PROCESSING INSTRUCTION (Conclusions, 4 to 6 paragraphs): Provide a clear
overview of what was accomplished across the four simulations and this
paper. Structure the section as follows:

(1) One paragraph that lists the four simulations and the headline
artifact counts (Simulation 1: 84 hours x 7 files = 588 files at minute
resolution across 4 sites and 116 robots; Simulation 2: 10 stage Python
scripts plus 30 ASCII progress diagrams plus 6 deliverable diagrams plus
6 regulatory tables for one patient (PAT-2026-0042) over 1,120 days;
Simulation 3: 24 hourly Python scripts, 24 JSON outputs, 75 ASCII
diagrams, 53 core agents organized into 4 layers; Simulation 4: 168
hourly Python scripts, 168 JSON outputs, 525 ASCII diagrams, 7 daily
summaries, 7 branches, 168 GitHub commits, 7 PRs, plus a complete
local-verification run on a 2015 Core i5-6200U laptop with 4 GB RAM).

(2) One paragraph on the persistent themes: repository-scale LLM context
(1M tokens) as the computational signature that distinguishes this work
from current narrow supervised oncology models; hourly commit cadence as
the operational signature aligned with the FDA RTCT vision; ASCII
diagrams plus markdown plus Python plus JSON as the file-level artifacts
that make the work auditable and human-monitorable in real time.

(3) One paragraph on the implications for patient prediction safety and
effectiveness: LLM code-and-text contextual awareness at scale enables
the best opportunity to improve patient prediction safety and effectiveness
in oncology trials, because (a) the model can reason across protocol,
EHR, robotic telemetry, regulatory adaptations, and ASCII diagrams in a
single inference pass, (b) the model can emit structured artifacts (JSON,
Python, ASCII) that downstream systems can ingest, and (c) the model's
output can be committed to GitHub at a 1-hour cadence, mirroring the
real-time signal-sharing framework the FDA validated through Paradigm
Health. Frame this implication respectfully with respect to the FDA and
to current trial sponsors; the four simulations are opportunistic
extensions of the FDA RTCT vision, not critiques of it.

(4) One paragraph that restates the key limitations: no real-patient
testing, non-deterministic LLM output, partial local verification on
constrained hardware, and the gap to TRIPOD+AI and CREMLS reporting
requirements. State that the work is a computational and software
demonstration, not a clinical demonstration.

(5) One paragraph that points forward: the next Claude Code 4.7 Max
generation pass will populate this template into a full 70+ page paper
with color diagrams, the author will then perform the senior-author
white-space and table-width formatting pass, and the artifact at
new-trial/national-24-7-trial/paper/LaTeX\_Source\_Files.zip will be
ready for upload to Overleaf. Cite \cite{fda2026realtime},
\cite{kawchak\_2026\_19176370}, \cite{kawchak\_2026\_19119939},
\cite{kawchak\_2026\_19396256}, \cite{Collins2024TRIPODAI}, and
\cite{ElEmam2024CREMLS}.]
